\documentclass[aspectratio=169,11pt]{beamer}
\usetheme{CambridgeUS}
\usecolortheme{default}

\definecolor{sapphire}{rgb}{0.03,0.15,0.4}
\definecolor{nupurple}{RGB}{78,42,132}
\definecolor{nugray}{RGB}{99,102,106}
\definecolor{orange}{rgb}{1,0.75,0}
\definecolor{almond}{rgb}{0.94,0.87,0.8}
\setbeamercolor{title}{fg=nupurple}
\setbeamercolor{frametitle}{fg=nupurple}
\setbeamercolor{structure}{fg=nupurple}
\setbeamercolor{section in head/foot}{bg=nupurple,fg=white}
\setbeamercolor{subsection in head/foot}{bg=nugray,fg=white}
\setbeamertemplate{navigation symbols}{}
\setbeamertemplate{itemize items}[circle]
\setbeamertemplate{itemize subitem}[triangle]
\setbeamertemplate{blocks}[rounded][shadow=true]
\setbeamercolor{block title}{fg=white,bg=nupurple}
\setbeamercolor{block body}{bg=nupurple!7,fg=black}
\setbeamercolor{alertblock title}{fg=white,bg=orange!55!black}
\setbeamercolor{alertblock body}{bg=orange!12,fg=black}
\setbeamercolor{exampleblock title}{fg=white,bg=teal!65!black}
\setbeamercolor{exampleblock body}{bg=teal!8,fg=black}
\setbeamersize{text margin left=0.65cm,text margin right=0.65cm}
\setlength{\tabcolsep}{5pt}

\usepackage[T1]{fontenc}
\usepackage[utf8]{inputenc}
\usepackage{lmodern}
\usepackage{ragged2e}
\usepackage{amsmath,amsfonts,amssymb,mathtools,bm}
\usepackage{graphicx,booktabs,tabularx,multicol,array,float}
\usepackage{hyperref,pdfpages,xcolor,tcolorbox,tikz}
\usetikzlibrary{shapes,fit,arrows.meta,positioning}
\newcommand{\E}{\mathbb{E}}
\newcommand{\Var}{\operatorname{Var}}
\newcommand{\R}{\mathbb{R}}
\newcommand{\ind}{\mathbf{1}}
\newcommand{\key}[1]{\textbf{\color{nupurple}#1}}
\newcommand{\term}[1]{\textbf{\color{nupurple}#1}}
\newcommand{\logit}{\operatorname{logit}}

\title[A Guide to Statistical Network Modeling -- Session 1.1]{A Guide to Statistical Network Modeling}
\subtitle{Session 1.1: Foundations of Static Exponential-Family Random Graph Models}
\author[Moses Boudourides]{%
\begin{tabular}{c}
{\color{nupurple}\Large Moses Boudourides}\\[0.6em]
{\normalsize \textit{Faculty, Graduate Program on Data Science}}\\
{\normalsize \textit{Northwestern University}}\\[0.6em]
\href{mailto:Moses.Boudourides@northwestern.edu}{\textcolor{teal}{\small\texttt{Moses.Boudourides@northwestern.edu}}}
\end{tabular}}
\date[October 28, 2026]{\begin{tabular}{c}
\Large{\href{https://instats.org/}{\texttt{\textcolor{teal}{instats}}} Seminar}\\[0.5em]
\small{Wednesday, October 28, 2026 \quad 3:00--4:30 PM GMT (90 minutes)}
\end{tabular}}

\begin{document}
\includepdf[pages={1}]{first_slide.pdf}
\includepdf[pages={2}]{second_slide.pdf}
\begin{frame}[plain]\titlepage\end{frame}

\begin{frame}{Companion computations and reproducibility}
\begin{itemize}
\item \term{Public repository:} \href{https://github.com/mboudour/no-code-statistical-network-modeling}{\textcolor{blue}{\underline{github.com/mboudour/no-code-statistical-network-modeling}}}.
\item \term{Companion Streamlit app:} \href{https://no-code-statistical-network-modeling.onrender.com/}{\textcolor{blue}{\underline{no-code-statistical-network-modeling.onrender.com}}}; open \term{Five worked public networks} for the five guided public-data workflows.
\item \term{No-code workflow:} select a documented public network, inspect its support and covariates, specify a static ERGM, fit it with the standard \texttt{statnet} engine, and preserve the output record.
\item \term{Scope of this session:} binary \emph{static} graph distributions; curved terms, degeneracy-aware refinement, and the full standard ERGM diagnostic audit are developed in Session 1.2, while temporal models are developed later.
\item \term{Principle:} an application can make computations reproducible; it cannot determine a credible network boundary, causal interpretation, or model specification automatically.
\end{itemize}
\end{frame}

\begin{frame}{Learning objectives}
By the end of this session, we will be able to:
\begin{itemize}
\item define a network observation, its \term{support}, and its \term{risk set};
\item distinguish \term{edge independence}, \term{dyad independence}, and structural dependence;
\item state the canonical static \term{exponential-family random graph model} (ERGM);
\item use \term{change statistics} to interpret conditional log-odds;
\item translate substantive network questions into a modest, documented ERGM specification; and
\item identify claims that an ERGM output does \emph{not} establish.
\end{itemize}
\end{frame}

\begin{frame}{Session route}
\begin{enumerate}
\item \term{Data and support:} what graph is observed, and which dyads are possible?
\item \term{Dependence:} why independent-dyad regression is often inadequate.
\item \term{Static ERGM:} probability law, statistics, parameters, and conditional interpretation.
\item \term{Specification:} structural and covariate statistics chosen from a research question.
\item \term{Computation:} likelihood, simulation, pseudolikelihood, and five public-network workflows.
\end{enumerate}
\end{frame}

\begin{frame}{Why statistical network modeling?}
\begin{itemize}
\item Conventional regression treats observations as independent conditional on covariates; a \term{network tie} can depend on other ties.
\item \term{Reciprocity:} $Y_{ij}$ may be informative about $Y_{ji}$.
\item \term{Closure:} a tie from $i$ to $j$ can be associated with paths through third actors.
\item \term{Degree heterogeneity:} ties sharing an actor can reflect common activity or popularity processes.
\item A statistical network model therefore defines a probability distribution for an \term{entire graph}, not an unrelated collection of rows.
\end{itemize}
\end{frame}

\begin{frame}{Network data: actors, relations, and graphs}
\begin{itemize}
\item Let $V=\{1,\ldots,n\}$ be the \term{actor set}; actors can be people, organizations, places, species, or events.
\item Let $Y_{ij}$ record the relation from $i$ to $j$; for a binary tie, $Y_{ij}\in\{0,1\}$.
\item The adjacency matrix $Y=(Y_{ij})$ is random; the observed network is its realization $y$.
\item The \term{unit of probabilistic modeling} is $Y$, while an actor, dyad, or local configuration can enter a statistic of $Y$.
\item A network question must state the relation, observation window, actor boundary, and whether a tie is directed or valued.
\end{itemize}
\end{frame}

\begin{frame}{Directed, undirected, and loopless support}
\begin{itemize}
\item In a \term{directed} network, $Y_{ij}$ and $Y_{ji}$ are distinct potential ties.
\item In an \term{undirected} network, $Y_{ij}=Y_{ji}$ and an unordered pair $\{i,j\}$ has one edge variable.
\item A graph is \term{loopless} when $Y_{ii}=0$; this is common but must be stated rather than assumed universally.
\item For a directed loopless binary graph, there are $n(n-1)$ potential ordered dyads.
\item For an undirected loopless binary graph, there are $\binom{n}{2}=n(n-1)/2$ potential unordered dyads.
\end{itemize}
\end{frame}

\begin{frame}{One-mode and two-mode networks}
\begin{itemize}
\item A \term{one-mode} network has ties among actors in one actor set, such as firms linked to firms.
\item A \term{two-mode} or \term{bipartite} network has two disjoint modes, such as women connected to events.
\item Bipartite support admits ties only across modes; within-mode dyads are \term{structural zeros}, not observed absences.
\item Projecting a bipartite graph to one mode changes the data object, induces dependence, and must not be treated as an innocuous formatting step.
\item An ERGM specification must use statistics and support appropriate to the network mode structure.
\end{itemize}
\end{frame}

\begin{frame}{Static is an observation design, not a claim of no change}
\begin{itemize}
\item A \term{static ERGM} models the distribution of a network observed at one designated time or aggregation window.
\item The network may have evolved before observation; a static model does not identify its unobserved sequence of tie changes.
\item Aggregating contacts over a long window can convert repeated events into a binary ``ever tied'' relation and can conceal duration or ordering.
\item A static interpretation is appropriate only for a question about the graph distribution at the stated observation design.
\item Network evolution is treated with TERGMs, STERGMs, and SAOMs in later sessions.
\end{itemize}
\end{frame}

\begin{frame}{The admissible dyad set and graph support}
For a loopless graph, define the admissible dyad set
\[
D=\begin{cases}
\{(i,j):i\ne j\}, & \text{directed},\\
\{(i,j):i<j\}, & \text{undirected}.
\end{cases}
\]
\begin{itemize}
\item The \term{sample space} $\mathcal{Y}$ is the set of all graph configurations allowed by the design.
\item With no additional constraints, a binary undirected graph has $|\mathcal{Y}|=2^{\binom n2}$ configurations.
\item Structural constraints, bipartite mode membership, and forbidden relations alter $\mathcal{Y}$ and therefore alter the probability model.
\end{itemize}
\end{frame}

\begin{frame}{Network boundary and risk set}
\begin{itemize}
\item The \term{boundary} specifies which actors are included in $V$.
\item The \term{risk set} specifies which dyads could in principle form or be observed during the study design.
\item Excluding actors who carry relational traffic can distort observed degree, mixing, and closure configurations.
\item An ERGM parameter is conditional on the selected actor set and support; it is not automatically portable to an unobserved wider network.
\item Report how actors entered the network and which possible dyads were excluded by design.
\end{itemize}
\end{frame}

\begin{frame}{Observed zero, missing tie, and structural zero}
\begin{itemize}
\item An \term{observed zero} is an admissible dyad recorded as absent.
\item A \term{missing tie} is an admissible dyad whose state is unobserved; it requires an observation-process assumption or an observed-data treatment.
\item A \term{structural zero} is a dyad not in the support, such as a within-mode pair in a bipartite network.
\item Coding missing dyads as zeros can manufacture sparsity and bias structural statistics.
\item Removing structural zeros after fitting is also incorrect: support must be specified before the likelihood is defined.
\end{itemize}
\end{frame}

\begin{frame}{Actor and dyadic covariates}
\begin{itemize}
\item An \term{actor covariate} $x_i$ is measured on an actor: role, seniority, wealth, group, or behavior.
\item A \term{dyadic covariate} $x_{ij}$ is measured on a pair: distance, co-membership, exposure, or opportunity.
\item Covariates must be aligned to the actor identities and observation design that define $Y$.
\item A covariate measured after the observed relation can be useful descriptively but may be unsuitable as a causal predictor.
\item ERGMs represent covariates through specified network statistics rather than through a generic row-level regression design matrix.
\end{itemize}
\end{frame}

\begin{frame}{Descriptive statistics versus generative parameters}
\begin{itemize}
\item A \term{descriptive statistic} summarizes the observed graph, for example edge count or clustering coefficient.
\item A \term{generative statistic} $g_k(Y,X)$ enters a probability model for many possible graphs in $\mathcal{Y}$.
\item A triangle count observed in $y$ does not by itself imply a positive ERGM closure parameter.
\item A homophilous observed mixing pattern can reflect group composition, degree heterogeneity, reciprocity, or other modeled features.
\item ERGM interpretation is therefore \term{conditional on the full specification}, not a one-statistic narrative.
\end{itemize}
\end{frame}

\begin{frame}{Relational dependence}
\begin{itemize}
\item \term{Relational dependence} means that the distribution of one tie can depend on other ties after covariates are considered.
\item Dependence is not an unspecified residual-correlation adjustment; it is represented through declared graph configurations.
\item A reciprocity statistic specifies a within-dyad dependence mechanism.
\item A closure statistic specifies dependence involving triples or longer local configurations.
\item The substantive task is to choose statistics that encode a plausible relational process while avoiding indiscriminate term accumulation.
\end{itemize}
\end{frame}

\begin{frame}{Edge independence and dyad independence}
\begin{itemize}
\item \term{Edge independence} assumes all admissible edge variables are mutually independent conditional on covariates.
\item In an undirected binary graph, edge independence and dyad independence coincide because each unordered dyad contains one edge variable.
\item In a directed graph, \term{dyad independence} allows $(Y_{ij},Y_{ji})$ to be dependent while distinct unordered dyads remain independent.
\item A reciprocity ERGM can be dyad independent but not edge independent.
\item This distinction determines which factorization and conditional interpretation are available.
\end{itemize}
\end{frame}

\begin{frame}{A local-dependence illustration}
\begin{itemize}
\item Under the Frank--Strauss Markov condition for an undirected graph, disjoint-edge variables satisfy
\[
Y_{ij}\ \perp\ Y_{kl}\mid Y_{-(ij),(kl)}\qquad\text{when }\{i,j\}\cap\{k,l\}=\varnothing.
\]
\item A positive distribution satisfying this condition can be represented using edge, $k$-star, and triangle configurations.
\item General ERGMs need not impose this particular Markov restriction; they may represent broader or nonlocal dependence.
\item The selected sufficient statistics, not the acronym ``ERGM,'' carry the substantive dependence assumptions.
\end{itemize}
\end{frame}

\begin{frame}{The canonical static ERGM}
\[
\Pr_{\boldsymbol\theta}(Y=y\mid X)
=\frac{\exp\{\boldsymbol\theta^{\mathsf T}g(y,X)\}}
{\kappa(\boldsymbol\theta,X)},
\qquad
\kappa(\boldsymbol\theta,X)=\sum_{y'\in\mathcal{Y}}
\exp\{\boldsymbol\theta^{\mathsf T}g(y',X)\}.
\]
\begin{itemize}
\item $g(y,X)=(g_1,\ldots,g_p)^{\mathsf T}$ is the vector of specified \term{network statistics}.
\item $\boldsymbol\theta\in\R^p$ is the canonical parameter vector.
\item $\kappa$ is the \term{normalizing constant}; it assigns total probability one over the selected support.
\end{itemize}
\end{frame}

\begin{frame}{Why the normalizing constant matters}
\begin{itemize}
\item $\kappa(\boldsymbol\theta,X)$ sums over every graph allowed by $\mathcal{Y}$.
\item Its size grows exponentially in the number of possible dyads; direct enumeration is infeasible except for very small graphs.
\item The intractable normalizing constant motivates simulation-based estimation and simulation-based diagnostics.
\item It does \emph{not} disappear from a likelihood merely because the observed graph is represented by an adjacency matrix.
\item Some comparative expressions are tractable because $\kappa$ cancels, as in graph probability ratios and change-statistic interpretations.
\end{itemize}
\end{frame}

\begin{frame}{Exponential-family identities}
Let $\psi(\boldsymbol\theta;X)=\log\kappa(\boldsymbol\theta,X)$. Then, for the canonical family,
\[
\nabla_{\boldsymbol\theta}\psi(\boldsymbol\theta;X)=\E_{\boldsymbol\theta}\{g(Y,X)\mid X\},
\qquad
\nabla^2_{\boldsymbol\theta}\psi(\boldsymbol\theta;X)=\Var_{\boldsymbol\theta}\{g(Y,X)\mid X\}.
\]
\begin{itemize}
\item The likelihood score compares observed sufficient statistics with their model expectations.
\item The second identity links curvature with variability of the statistics and expected Fisher information.
\item These canonical identities require modification for curved parameterizations, treated in Session 1.2.
\end{itemize}
\end{frame}

\begin{frame}{Probability ratios compare complete graphs}
For $y_a,y_b\in\mathcal{Y}$ under the same covariates,
\[
\frac{\Pr_{\boldsymbol\theta}(Y=y_a\mid X)}
{\Pr_{\boldsymbol\theta}(Y=y_b\mid X)}
=\exp\!\left[\boldsymbol\theta^{\mathsf T}\{g(y_a,X)-g(y_b,X)\}\right].
\]
\begin{itemize}
\item The normalizing constant cancels in the ratio.
\item Parameters assign relative weight to differences in network statistics.
\item This is a comparison of \term{entire network configurations}, not a marginal probability for a randomly chosen tie.
\item A one-dyad toggle yields the conditional-logit interpretation on the next slides.
\end{itemize}
\end{frame}

\begin{frame}{Change statistics}
For an admissible binary dyad $(i,j)$, let $y^+_{ij}$ and $y^-_{ij}$ agree except that the dyad is set to one or zero. Define
\[
\Delta_{ij}g(y,X)=g(y^+_{ij},X)-g(y^-_{ij},X).
\]
\begin{itemize}
\item A \term{change statistic} records how each modeled feature would change if one tie were toggled while all other ties are held fixed.
\item It depends on $y_{-ij}$ and $X$, but not on the current value of $Y_{ij}$.
\item Change statistics make local structural context explicit: the same potential tie can have different model contribution in different network neighborhoods.
\end{itemize}
\end{frame}

\begin{frame}{Conditional-logit identity}
\[
\logit\Pr_{\boldsymbol\theta}(Y_{ij}=1\mid Y_{-ij}=y_{-ij},X)
=\boldsymbol\theta^{\mathsf T}\Delta_{ij}g(y,X).
\]
\begin{itemize}
\item Each coefficient weights the corresponding \term{change statistic} in the conditional log-odds of a tie.
\item ``Conditional'' means conditional on the complete remainder of the graph and included covariates.
\item For a toggle with $\Delta_{ij}g_k=1$, exponentiating $\theta_k$ gives the multiplicative contribution of that term to the conditional odds, holding the other change statistics fixed.
\item This is not a population-average tie probability and is not automatically a causal effect.
\end{itemize}
\end{frame}

\begin{frame}{Toggle interpretation requires support}
\begin{itemize}
\item The conditional-logit identity applies when both $y^+_{ij}$ and $y^-_{ij}$ belong to $\mathcal{Y}$.
\item A structural zero has no binary toggle comparison because its tie-present state is not admissible.
\item Constrained supports can restrict which single-edge toggles are valid; for example, fixed-degree support generally requires degree-preserving moves rather than arbitrary single-edge toggles.
\item Bipartite support permits toggles only across modes.
\item Before interpreting a coefficient, verify that the corresponding tie comparison is defined under the network's stated support.
\item Software options for directedness, bipartite structure, missing ties, and constraints are statistical assumptions, not merely formatting controls.
\end{itemize}
\end{frame}

\begin{frame}{The edges statistic: baseline sparsity}
\begin{itemize}
\item For a binary graph, $g_{\mathrm{edges}}(y)=\sum_{(i,j)\in D}y_{ij}$ counts present ties.
\item Its coefficient controls baseline conditional tie propensity after all other specified effects are included.
\item A negative edges coefficient is common in sparse networks, but its magnitude depends on the support and the rest of the model.
\item It is not simply the observed density transformed by a logistic function once structural and covariate terms are included.
\item Every introductory ERGM needs a clearly interpreted baseline relation term.
\end{itemize}
\end{frame}

\begin{frame}{Reciprocity in a directed network}
\[
g_{\mathrm{mutual}}(y)=\sum_{i<j}y_{ij}y_{ji}.
\]
\begin{itemize}
\item The statistic counts mutually tied unordered pairs in a directed graph.
\item Its change statistic is one when toggling $i\to j$ completes a reciprocal dyad, and zero otherwise.
\item A positive coefficient favors reciprocated dyads \emph{relative to} otherwise comparable directed graphs under the full model.
\item Reciprocity is not defined for an undirected graph and must not be added mechanically to a one-mode undirected or bipartite specification.
\end{itemize}
\end{frame}

\begin{frame}{Degree patterns and centralization}
\begin{itemize}
\item The \term{degree} of an undirected actor counts incident ties; directed graphs distinguish indegree and outdegree.
\item A degree distribution is descriptive; it can reveal isolates, hubs, and dispersion that a simple edges-only model will not reproduce.
\item Star statistics count configurations sharing a focal actor and can represent degree-related structure.
\item Raw high degree is not evidence that a particular star parameter belongs in a model; the statistic must answer a substantive question and be diagnostically stable.
\item Curved geometrically weighted degree terms are deferred to Session 1.2.
\end{itemize}
\end{frame}

\begin{frame}{Closure and triadic structure}
\begin{itemize}
\item A triangle in an undirected graph is a fully connected triple; directed closure has several distinct configurations.
\item Closure terms can represent a tendency for ties to occur in locally connected neighborhoods.
\item Triangle counts overlap strongly: a single tie can complete several triangles, so the change statistic depends on shared partners.
\item A positive triangle-related parameter is conditional on the other statistics; it is not proof of a social mechanism such as trust or contagion.
\item Session 1.2 treats stable curved alternatives to naïve high-order structural terms.
\end{itemize}
\end{frame}

\begin{frame}{Actor-level covariate effects}
For a scalar attribute $x_i$ in a directed graph, possible statistics include
\[
g_{\mathrm{send}}(y,x)=\sum_{i\ne j}y_{ij}x_i,
\qquad
g_{\mathrm{recv}}(y,x)=\sum_{i\ne j}y_{ij}x_j.
\]
\begin{itemize}
\item A \term{sender} effect asks whether actors with larger $x_i$ tend to send more ties.
\item A \term{receiver} effect asks whether actors with larger $x_j$ tend to receive more ties.
\item These are distinct hypotheses; a single correlation between degree and an attribute cannot distinguish them.
\item Interpret both conditionally on structural and other covariate terms included in the model.
\end{itemize}
\end{frame}

\begin{frame}{Homophily and mixing}
For a categorical attribute $x_i$, an undirected matching statistic is
\[
g_{\mathrm{match}}(y,x)=\sum_{i<j}y_{ij}\ind(x_i=x_j).
\]
\begin{itemize}
\item A positive matching coefficient favors same-category ties relative to cross-category ties after other terms are held fixed.
\item A \term{mixing matrix} represents ties through category-pair statistics, allowing category-specific patterns that a single matching statistic can conceal.
\item Directed and undirected mixing statistics must be defined separately; an undirected category-pair representation must not double-count dyads.
\item Homophily coefficients do not automatically identify preference: opportunity, group size, unmeasured similarity, and boundary design may contribute.
\end{itemize}
\end{frame}

\begin{frame}{A coefficient is conditional, not causal}
\begin{itemize}
\item ERGM parameters describe relative probability contributions conditional on the rest of the graph, selected support, covariates, and model terms.
\item A positive covariate or homophily effect is not by itself an intervention effect.
\item A positive closure term is not direct evidence of a psychological preference for triangles.
\item Causal claims require a design, temporal ordering, treatment definition, and explicit account of confounding and interference.
\item Scientific interpretation should separate \term{model-based association} from a stronger causal explanation.
\end{itemize}
\end{frame}

\begin{frame}{From research question to ERGM specification}
\begin{enumerate}
\item Define actors, relation, time window, boundary, and support.
\item State a graph-level question: reciprocity, mixing, closure, degree concentration, or covariate-associated ties.
\item Select a small set of statistics justified by the question and network type.
\item Decide which covariates are pre-relation, observed, and aligned to actors or dyads.
\item Fit, inspect computational diagnostics, compare simulated and observed features, and revise only with a documented rationale.
\end{enumerate}
\end{frame}

\begin{frame}{Likelihood and the estimation problem}
\begin{itemize}
\item For observed graph $y$, the log likelihood is
\[
\ell(\boldsymbol\theta)=\boldsymbol\theta^{\mathsf T}g(y,X)-\log\kappa(\boldsymbol\theta,X).
\]
\item The score equates observed statistics with their model expectations.
\item Direct maximum likelihood requires expectations under a distribution with an intractable normalizing constant.
\item \term{Monte Carlo maximum likelihood estimation} approximates those expectations by simulating networks from trial parameter values.
\item Convergence, sample-space behavior, and simulation quality are part of the statistical result, not post-processing details.
\end{itemize}
\end{frame}

\begin{frame}{MCMC maximum likelihood: conceptual cycle}
\begin{enumerate}
\item Start from a parameter value and simulate networks by admissible tie toggles or other valid moves.
\item Estimate expected sufficient statistics under the current model.
\item Compare simulation expectations with $g(y,X)$ and update the parameter value.
\item Repeat until the estimating discrepancy is sufficiently small under stated convergence criteria.
\item Retain settings, random seed policy, warnings, and diagnostics in the reproducibility record.
\end{enumerate}
\end{frame}

\begin{frame}{Maximum pseudolikelihood: useful but limited}
\begin{itemize}
\item \term{Maximum pseudolikelihood estimation} (MPLE) forms a pseudolikelihood by multiplying the full conditional probabilities of individual dyads, rather than evaluating the joint ERGM likelihood.
\item It is computationally convenient and can provide a starting value for MCMC estimation.
\item Under dependence, this product of full conditionals is not the joint ERGM likelihood; standard errors and point estimates can be biased or overconfident relative to the likelihood target.
\item A logistic-regression-looking output is not evidence that dependence has disappeared.
\item The app labels pseudolikelihood as an approximation; reported primary results use the standard simulation-based ERGM engine when it converges.
\end{itemize}
\end{frame}

\begin{frame}{Uncertainty and inferential discipline}
\begin{itemize}
\item A standard error quantifies uncertainty conditional on the fitted model, data boundary, support, and estimation procedure.
\item A confidence interval does not repair an implausible statistic, a poor network boundary, missing data, or failed MCMC convergence.
\item Model selection based only on coefficient significance risks specification search and selective reporting.
\item Report estimates, uncertainty, simulation settings, diagnostics, and the scope of the population to which the observed network pertains.
\item Avoid dichotomous conclusions such as ``proved a network mechanism.''
\end{itemize}
\end{frame}

\begin{frame}{ERGM diagnostic audit}
\begin{itemize}
\item \term{Estimation / MCMC:} inspect trace plots, autocorrelation plots, and sampled-statistic distributions.
\item \term{Goodness of fit:} compare observed and simulated degree, geodesic-distance, edgewise shared-partner, and dyadwise shared-partner distributions when they are meaningful for the support.
\item \term{Model-specific checks:} examine reciprocity, mixing, indegree/outdegree, triad census, component size, or other substantively important features not directly fitted.
\item A converged MCMC chain does not imply good model fit: convergence concerns estimation behavior, whereas goodness of fit concerns the distribution of networks generated by the fitted model.
\item Retain diagnostic plots, simulation settings, and unresolved discrepancies in the reproducibility record; the full standard diagnostic workflow is developed in Session 1.2.
\end{itemize}
\end{frame}

\begin{frame}{Five worked public networks in the companion app}
\small
\begin{tabularx}{\linewidth}{@{}lXl@{}}
\toprule
\term{Network} & \term{Network type and introductory workflow} & \term{Static support}\\
\midrule
Florentine families & Business relations; edges and documented family attributes & Undirected, binary, one-mode\\
Sampson's monastery & Interpersonal nominations; edges, reciprocity, group mixing & Directed, binary, one-mode\\
Lazega law firm & Professional advice; edges, reciprocity, office or status covariates & Directed, binary, one-mode\\
Southern Women & Woman--event participation; bipartite edges and mode structure & Undirected, binary, two-mode\\
Kapferer tailor shop & Frequent sociational relations; density and documented closure comparison & Undirected, binary, one-mode\\
\bottomrule
\end{tabularx}
\vspace{3pt}
Each example is accompanied by source, data dictionary, validity warning, and reproducible computation record.
\end{frame}

\begin{frame}{Worked example: Florentine families}
\begin{itemize}
\item \term{Question:} are business ties associated with documented family attributes after accounting for baseline density?
\item \term{Data object:} one static, undirected, binary family business network with actor attributes.
\item \term{Foundational specification:} an edges term plus a pre-specified attribute term or matching statistic that is meaningful for the recorded families.
\item \term{Computation:} inspect estimation diagnostics and conditional coefficient estimates, then compare substantively relevant observed network features with their distributions under simulated networks.
\item \term{Limit:} the small historical network supports a focused illustration, not a general causal claim about wealth or political influence.
\end{itemize}
\end{frame}

\begin{frame}{Worked example: Sampson's monastery}
\begin{itemize}
\item \term{Question:} do interpersonal nominations exhibit reciprocity and within-group mixing?
\item \term{Data object:} a directed, binary, one-mode sociometric network among monks with recorded group affiliation.
\item \term{Foundational specification:} edges + mutual + group matching, with terms included because each maps to a separate question.
\item \term{Interpretation:} the mutual coefficient concerns conditional odds of reciprocated nominations, not a marginal probability of friendship.
\item \term{Limit:} a single observation wave cannot establish the sequence by which affiliations and nominations co-developed.
\end{itemize}
\end{frame}

\begin{frame}{Worked example: Lazega law firm}
\begin{itemize}
\item \term{Question:} how are professional advice ties conditionally associated with reciprocity and documented organizational attributes?
\item \term{Data object:} a directed, binary, one-mode advice network with actor-level law-firm variables.
\item \term{Foundational specification:} start with edges and reciprocity, then add one documented attribute hypothesis rather than all available covariates.
\item \term{Computation:} retain the network boundary, directedness, and attribute coding in the app's analysis record.
\item \term{Limit:} a cross-sectional model describes the observed advice-network distribution; it does not identify influence or organizational change.
\end{itemize}
\end{frame}

\begin{frame}{Worked example: Davis Southern Women}
\begin{itemize}
\item \term{Question:} how does an affiliation graph differ from a one-mode social graph?
\item \term{Data object:} women linked to social events; a static, binary, bipartite incidence network.
\item \term{Support:} only woman--event ties are admissible; woman--woman and event--event cells are structural zeros in the original graph.
\item \term{Computation:} inspect mode sizes and degree distributions before fitting a bipartite-appropriate baseline ERGM.
\item \term{Limit:} a one-mode projection creates a different network and requires a separately justified analysis.
\end{itemize}
\end{frame}

\begin{frame}{Worked example: Kapferer tailor shop}
\begin{itemize}
\item \term{Question:} do relatively frequent sociational relations display local closure after the baseline density is represented?
\item \term{Data object:} a static, undirected, binary, one-mode network of 39 tailor-shop workers from the first collection period.
\item \term{Foundational specification:} begin with an edges-only model; treat a closure extension as a documented later comparison, not an automatic term.
\item \term{Computation:} inspect support, density, degree distribution, MCMC behavior, and simulated graph summaries before interpreting an added structural effect.
\item \term{Limit:} the 39-worker object is a selected analytic panel and records thresholded frequent interaction, not all social contact in a probability sample of factories.
\end{itemize}
\end{frame}

\begin{frame}{No-code Session 1.1 computation workflow}
\begin{enumerate}
\item Select one of the five public networks and read its source, network-type, and scope note.
\item Inspect actor count, admissible dyads, density, directionality, missingness, attributes, and degree distribution.
\item Choose a small supported formula from the session's available terms; the app prevents directed-only terms in undirected data.
\item Run the standard \texttt{statnet} ERGM computation; examine the coefficient table, MCMC messages, and model limitations before the full diagnostic audit in Session 1.2.
\item Download a Markdown computation record containing data source, support, formula, software settings, results, warnings, and limitations.
\end{enumerate}
\end{frame}

\begin{frame}{BYOD: required upload contract}
\begin{itemize}
\item \term{Edge table:} source ID and target ID for observed binary ties; a valued relation is not silently dichotomized and belongs in an explicitly justified later workflow.
\item \term{Node table:} one unique actor ID per row and optional pre-specified actor attributes.
\item \term{Design choices:} directed versus undirected; one-mode versus bipartite; actor boundary; structural zeros; missingness; and time window.
\item \term{Minimum checks:} duplicate dyads, self-loops, unmatched node IDs, empty attributes, implausible density, and unsupported formula terms.
\item \term{Principle:} upload does not make an ERGM appropriate; the research question and support must justify the model.
\end{itemize}
\end{frame}

\begin{frame}{Common invalid or misleading uses}
\begin{itemize}
\item Fitting an undirected model to a directed relation after discarding direction without a substantive reason.
\item Treating unobserved dyads as zeros or ignoring an actor boundary that omits important relational alternatives.
\item Adding every available structural statistic or covariate merely because software accepts the formula.
\item Interpreting a coefficient as a marginal probability, an intervention effect, or a proof of a social mechanism.
\item Reporting a pseudolikelihood result as if it had resolved the ERGM normalizing-constant problem.
\end{itemize}
\end{frame}

\begin{frame}{A reporting checklist for a static ERGM}
\begin{itemize}
\item \term{Data:} actors, relation, observation window, boundary, directionality, mode structure, missingness, and support constraints.
\item \term{Question:} which structural or covariate-associated graph feature is under investigation?
\item \term{Specification:} exact sufficient statistics, covariate coding, and rationale for each term.
\item \term{Estimation:} software, package version, algorithm, simulation controls, convergence status, and warnings.
\item \term{Results:} conditional coefficient interpretation, uncertainty, simulated comparisons, limitations, and scope of inference.
\end{itemize}
\end{frame}

\begin{frame}{Session 1.1 synthesis}
\begin{itemize}
\item A static ERGM is a probability model for a graph on an explicitly defined support.
\item Its parameters weight \term{network statistics}; they are interpreted through conditional changes under the rest of the graph.
\item Data definition, boundaries, structural zeros, and covariate timing are part of the statistical model.
\item For most nontrivial dependent ERGMs, simulation-based methods such as MCMC are used because the normalizing constant cannot be evaluated directly.
\item Next: \term{curved ERGMs, stable structural terms, fit, and degeneracy} in Session 1.2.
\end{itemize}
\end{frame}

\begin{frame}{References}
\footnotesize
\begin{thebibliography}{99}
\bibitem{boudourides} Boudourides, M. (2026). \emph{A Guide to Statistical Network Modeling: ERGMs, Temporal ERGMs, and SAOMs}. Seminar guide.
\bibitem{wasserman} Wasserman, S. and Faust, K. (1994). \emph{Social Network Analysis}. Cambridge University Press.
\bibitem{frank} Frank, O. and Strauss, D. (1986). Markov graphs. \emph{Journal of the American Statistical Association}, 81, 832--842.
\bibitem{robins} Robins, G., Pattison, P., Kalish, Y., and Lusher, D. (2007). An introduction to exponential random graph models for social networks. \emph{Social Networks}, 29, 173--191.
\bibitem{hunter} Hunter, D. R. and Handcock, M. S. (2006). Inference in curved exponential family models for networks. \emph{Journal of Computational and Graphical Statistics}, 15, 565--583.
\bibitem{handcock} Handcock, M. S. and Gile, K. J. (2010). Modeling social networks from sampled data. \emph{Annals of Applied Statistics}, 4, 5--25.
\end{thebibliography}
\end{frame}

\begin{frame}{Questions and discussion}
\begin{center}
\Large{Thank you!}\\[2ex]
\normalsize{Questions?}\\[5ex]
\href{mailto:Moses.Boudourides@northwestern.edu}{\textcolor{teal}{\small\texttt{Moses.Boudourides@northwestern.edu}}}\\[0.8em]
\href{mailto:Moses.Boudourides@gmail.com}{\textcolor{teal}{\small\texttt{Moses.Boudourides@gmail.com}}}
\end{center}
\end{frame}
\includepdf[pages={8}]{last_slide.pdf}
\end{document}
